\documentclass[11pt, reqno]{article}

\usepackage{amsmath, amsthm, amssymb}
\usepackage{enumitem}
\usepackage{tcolorbox}
\usepackage{hyperref}
\usepackage{tikz}
\usepackage{tikz-cd}
\usepackage{pgfplots}
\pgfplotsset{compat=1.18}
\usetikzlibrary{arrows.meta}
\usepackage{mathrsfs}
\usepackage{fancyhdr}
\usepackage[bottom=0.75in, top=1in, left=0.5in, right=0.5in]{geometry}
\usepackage{array}   % for \newcolumntype macro
\newcolumntype{L}{>{$}l<{$}}

\theoremstyle{plain}
\newtheorem*{theorem}{Theorem}
\newtheorem*{proposition}{Proposition}
\newtheorem{exercise}{Exercise}
\newtheorem*{lemma}{Lemma}
\newtheorem*{corollary}{Corollary}

\theoremstyle{definition}
\newtheorem*{definition}{Definition}
\newtheorem*{example}{Example}

\theoremstyle{remark}
\newtheorem*{remark}{Remark}

\renewcommand{\phi}{\varphi}
\renewcommand{\epsilon}{\varepsilon}
\renewcommand{\emptyset}{\varnothing}

\newcommand{\RR}{\mathbb{R}}
\newcommand{\ZZ}{\mathbb{Z}}
\newcommand{\NN}{\mathbb{N}}
\newcommand{\CC}{\mathbb{C}}
\newcommand{\QQ}{\mathbb{Q}}

\DeclareMathOperator{\ima}{\text{im}}

\begin{document}

\topmargin=-40pt
\rhead{Henry Woodburn}
\lhead{Math 622}
\renewcommand{\headrulewidth}{1pt}
\renewcommand{\headsep}{20pt}
\thispagestyle{fancy}

{\Huge \bfseries \noindent Homework 9}

\begin{enumerate}
    \item[1.] Let $M$ be a manifold and $\{U_i, \phi_i\}$ at atlas for $M$. Restricting to some neighborhood
    $U_i$, define a $(0,2)$ tensor $\omega_i = \phi_i^*(\cdot, \cdot)_{\RR^n}$, the pullback of the usual inner product 
    on $\RR^n$, where $n$ is the dimension of $M$. In order to show $\omega_i$ defines an inner product on every $T_p M$ 
    for every $p \in U_i$, we must show it is symmetric and postive definite. Fix $p \in U_i$. For any $u, v \in T_p M$, 
    \[
        \omega_i(u, v) = (d\phi(u), d\phi(v))_{\RR^n} = (d\phi(v), d\phi(u))_{\RR^n} = \omega_i(v,u),
    \]
    and 
    \[
        \omega_i(u,u) = (d\phi(u),d\phi(u))_{\RR^n} > 0
    \]
    unless $u = 0$ since $d\phi$ is injective. 

    Let $\{f_i\}$ be a partition of unity subordinate to $\{U_i\}$. We can extend $\omega_i$ to be a $(0,2)$-tensor 
    on all of $M$ by taking $\tilde{\omega}_i = f_i\omega_i$. Finally define a $(0,2)$-tensor 
    \begin{equation}\label{omega}
        \omega = \sum f_i \omega_i.
    \end{equation}
    I claim that $\omega$ defines a riemannian structure on $M$. The symmetry of $\omega$ is immediate from its definition.
    For positivity at each $T_p M$, suppose $\omega|_p(u,u) = 0$. Since each $\omega_i$ and $f_i$ is non-negative,
    every element of the sum \ref{omega} must be zero. In particular, if $U_j$ is an open set containing $p$,
    $\omega_j|_p(u,u) = 0$ and thus $u = 0$. 
    
    \item[2.] Let $F: U = \left(0,2\pi\right)\times \left(0,2\pi\right) \to T^2$ be the map 
    \[
        (s,t) \mapsto (cos(s), sin(s), cos(t), sin(t)).
    \]
    We view the domain $(0,2\pi)\times(0,2\pi)$ as an open submanifold of the quotient space 
    $\left([0,2\pi]\times[0,2\pi]\right)/\sim$, 
    which identifies each point on the boundary of the square with the one opposite to it. This space
    is diffeomorphic to our original $T^2$, since both are defined as products of copies of $S^1$. We compute
    the integral of the pullback $F^*\omega$ on the open submanifold, as $F$ is a diffeomorphism of the submanifold.
    Since the complement of the submanifold has measure zero, this integral will be equal to the integral
    over the entire manifold of the pullback. 

    By the formula for pullbacks, 
    \[
        F^*\omega = \cos(t)\sin^2(t)\sin^2(s).
    \]
    Then
    \[
        \int_T^2 \omega = \int_U F^* \omega = \int_0^{2\pi}\int_0^{2\pi}\cos(t)\sin^2(t)\sin^2(s)dt ds = 0,
    \]
    since the integral of cosine from 0 to $2\pi$ is zero, and we can separate variables in this integral. 

    \item[3.] (a.) Let $\alpha = \frac{1}{2\pi}\frac{x dy - y dx}{x^2 + y^2}$ be a 1-form on $\RR^2 \setminus \{0\}$.
    To show $\alpha$ is closed, we just calculate 
    \[
        d\alpha = \frac{x^2 + y^2 - 2x^2}{(x^2 + y^2)^2}dx \wedge dy - \frac{x^2 + y^2 - 2y^2}{(x^2 + y^2)^2}dy \wedge dx
        = 0.
    \]
    To integrate along the unit circle, map $\left[0,2\pi\right)$ to the unit circle by $t \mapsto (\cos(t),\sin(t))$.
    The pullback of $\alpha$ under this map is $\frac{1}{2\pi}(\cos^2(t) + \sin^2(t))dt = \frac{1}{2\pi}dt$. Then
    it is clear that the integral of $\alpha$ over $S^1$ is $1$. 

    By Stokes' theorem, the integral of any exact form over a manifold without boundary is zero. 
    Then $\alpha$ is not exact. Moreover, the integral of $\alpha$ over $S^1$ is equal to
    the integral of the pullback of $\alpha$ over $S^1$, since the pullback in this case is just the restriction
    to the circle. Then $i^*(\alpha)$ is not exact by the same reasoning.

    (b.) The cohomology group for top forms on $S^n$ is the one dimensional vector space $\RR$. The first 
    cohomology group of $\RR^2\setminus \{0\}$ is isomorphic to the first cohomology group of $S^1$ by
    homotopy invariance of de Rham cohomology. Then $H^1(\RR^n\setminus\{0\})$ is spanned by $\alpha$, and 
    we can choose $c \in \RR$ such that $\omega - c\alpha$ is zero in the quotient group of closed
    forms minus exact forms. In other words, $\omega - c\alpha$ is exact, and there exists a function $g$ 
    on $\RR^2 \setminus\{0\}$ such that $\omega = c\alpha + dg$.

    \item[4.] (a.) Let $\omega$ be a closed 1-form on $S^2$. Let $\gamma$ be any closed path in $S^2$. Then
    there is an open subset $U \subset S^2$ such that $\partial U = \gamma$ by the jordan curve theorem.
    Then by Stokes' theorem,
    \[
        \int_\gamma \omega = \int_{\partial U}\omega = \int_U d\omega = 0.
    \]

    Then the integral of $\omega$ over any closed curve is zero. This is equivalent to $\omega$ being exact. 

    (b.) We can calculate that
    \[
        d\omega = \frac{|x|^3 - 3x_1^2 |x| - 3x_2^2 |x| - 3x_3^2 |x|}{|x|^6} = 
        \frac{3|x|^3 - 3|x|^3}{|x|^6} = 0,
    \]
    thus $\omega$ is exact. 

    (c.) We make a coordinate change to spherical polar coordinates on $S^2$ and compute the integral
    of the pullback of $\omega$. 

    We use coordinates
    \begin{align*}
        & x = sin\theta \cos\phi\\
        & y = \sin\theta \sin\phi\\
        & z = \cos\phi.
    \end{align*}

    Let $T$ be the coordinate change. For the pullback, we calculate that 
    \begin{align*}
        & d(x(\theta, \phi)) = \cos\theta\cos\phi d\theta - \sin\theta\sin\phi d\phi\\
        & d(y(\theta, \phi)) = \cos\theta\sin\phi d\theta + \sin\theta\cos\phi d\phi\\
        & d(z(\theta, \phi)) = -\sin\theta d\theta.
    \end{align*}

    Then the pullback is 
    \[
        T^*\omega = (\sin^3 \theta + \cos^2\theta \sin\theta)d\theta \wedge d\phi = \sin\theta d\theta\wedge d\phi.
    \]

    and the integral is
    \[
        \int_S^2 \omega = \int_{(0,2\pi)\times (0,\pi)}T^*\omega = \int_0^{2\pi}\int_0^\pi \sin(\theta)d\theta d\phi
        = \int_0^{2\pi} 2 d\phi = 4\pi.
    \]

    This shows $\omega$ is not exact, because the integral of an exact form over a manifold without boundary
    is zero by Stokes' theorem.

    (d.) The hodge star map sends $p$ forms to $n-p$ forms with sign depending on the total orientation choice.
    Then we can see that 
    \[
        *\alpha = \sum_i (-1)^{i-1}\frac{x_i dx_1 \wedge \cdots \hat{dx_i}\wedge \cdots \wedge dx_n}{|x|^n},
    \]
    where the factor $(-1)^{i-1}$ is chosen so that the pairing of each $p-1$ form in the sum with its missing
    element $dx_i$ is correctly oriented after swapping $i-1$ times. 

    $*\alpha$ is a closed form, which can be shown similarly to previous items. We have to find the 
    differential of the coefficient $\frac{x_i}{|x|^n}$, which has the numerator $|x|^n - n |x|^(n-2)(x_i)$
    after permuting the $dx_i$ to the correct spot, which gets rid of the $(-1)^{i-1}$ coefficient. Once
    we sum over $i$, the terms will cancel just as in (b.). 

    (e.) $*\alpha$ is just the volume form for $S^n$. Then its integral over $S^n$ will equal 
    the surface area of the $n$-sphere. This proves $*\alpha$ is not exact by the reasoning above.

    \item[5.] (a.) Using Stokes' theorem and the product formula for the exterior derivative, we have
    \[
        0 = \int_{\partial M}\alpha \wedge \beta = \int_{M} d(\alpha \wedge \beta) = \int_M d\alpha \wedge \beta
        + (-1)^p\int_M \alpha \wedge d\beta.
    \]
    The result follows immediately after rearranging.

    (b.) It is proved in Lee that integration of top forms defines an isomorphism onto the top de Rham groups. 
    Since $\omega$ is exact, $\omega = d\alpha$ for some $\alpha$. Then $\omega \wedge \omega = d\alpha \wedge d\alpha
    = d(\alpha \wedge d\alpha)$, so $\omega \wedge \omega$ is exact. In particular, the integral over $M$
    must send it to zero. Then since the space of top forms is one dimensional, $\omega \wedge \omega$ is just 
    a smooth function. Then $\omega \wedge\omega$ is either positive and negative, and must equal zero somewhere
    by the intermediate value theorem, or is uniquely zero on $M$. 

    (c.) It is proved in Lee that the lower cohomology groups of $S^n$ are trivial, meaning the exact forms 
    are exactly the closed ones. Then if $\omega$ is closed, it is in fact exact, and the previous part applies.

    \item[6.] (a.) For symmetry, we have
    \[
        \{h_1, h_2\} = dh_2(\overrightarrow{h_1}) = i_{\overrightarrow{h_2}}\omega(h_1) = 
        \omega(\overrightarrow{h_1}, \overrightarrow{h_2}) = -\omega(\overrightarrow{h_2},\overrightarrow{h_1})
        = -\{h_2,h_1\}.
    \]

    (b.) Don't know unfortunately.

\end{enumerate}

\end{document}